% Part: first-order-logic
% Chapter: axiomatic-deduction
% Section: proving-things

\documentclass[../../../include/open-logic-section]{subfiles}

\begin{document}

\iftag{FOL}
      {\olfileid{fol}{axd}{pro}}
      {\olfileid{pl}{axd}{pro}}

\olsection{Conto \usetoken{P}{derivation}}


\begin{ex}
Upamana kita arep mbuktekake $(\lnot !D \lor !E) \lif (!D \lif
!E)$. Cetha yen iki dudu conto saka aksioma apa wae, mula kita kudu
nggunakake aturan \MP{} kanggo !!{derive} formula iku. Siji-sijine
aturan kita yaiku~MP, kang, yen diwenehi $!A$ lan $!A \lif !B$,
ngidini kita mbenerake~$!B$. Salah siji strategi yaiku nggunakake
\olref[prp]{ax:lor3} kanthi $!A$ minangka $\lnot !D$, $!B$ minangka
$!E$, lan $!C$ minangka $!D \lif !E$, yaiku conto
\[
(\lnot !D \lif (!D \lif !E)) \lif ((!E \lif (!D \lif !E)) \lif ((\lnot
!D \lor !E) \lif (!D \lif !E))).
\]
Apa alesane? Rong aplikasi MP ngasilake bagean pungkasan, yaiku formula
kang kita karepake. Lan kita gampang weruh yen $\lnot !D \lif (!D \lif
!E)$ iku conto saka \olref[prp]{ax:lnot2}, dene $!E \lif (!D \lif !E)$
iku conto saka \olref[prp]{ax:lif1}. Dadi, !!{derivation} kita yaiku:
\begin{derivation}
  1. &  $\lnot !D \lif (!D \lif !E)$ & \olref[prp]{ax:lnot2} \\
  2. & $(\lnot !D \lif (!D \lif !E)) \lif {}$\\
  &\qquad $((!E \lif (!D \lif !E)) \lif ((\lnot !D \lor !E) \lif (!D \lif !E)))$ & \olref[prp]{ax:lor3}\\
  3. & $(!E \lif (!D \lif !E)) \lif ((\lnot !D \lor !E) \lif (!D \lif !E))$ & 1, 2, \MP\\
  4. & $!E \lif (!D \lif !E)$ & \olref[prp]{ax:lif1}\\
  5. & $(\lnot !D \lor !E) \lif (!D \lif !E)$ & 3, 4, \MP
\end{derivation}
\end{ex}

\begin{ex}\ollabel{ex:identity}
Ayo kita nyoba golek !!a{derivation} saka $!D \lif !D$. Formula iku
dudu conto aksioma, mula kita kudu nggunakake \MP{} kanggo
!!{derive} formula kasebut. \olref[prp]{ax:lif1} iku aksioma kanthi
wujud~$!A \lif !B$ kang bisa dikenani~\MP. Mesthi wae, supaya migunani,
$!B$ kang dibenerake \MP{} minangka langkah kang bener ing kasus iki
kudu padha karo~$!D \lif !D$, amarga iki kang arep kita !!{derive}.
Iki ateges $!A$ uga kudu padha karo $!D$; mula kita bisa nimbang conto
\olref[prp]{ax:lif1} iki:
\[
!D \lif (!D \lif !D)
\]
Kanggo ngetrapake \MP, kita uga kudu mbenerake premis kapindho kang
cocog, yaiku~$!A$. Nanging ing kasus kita, premis iku bakal dadi~$!D$,
lan kita ora bakal bisa !!{derive}~$!D$ dhewe. Mula kita mbutuhake
strategi liyane.

Aksioma liyane kang mung ngemot~$\lif$ yaiku
\olref[prp]{ax:lif2}, yaiku
\[
(!A \lif (!B \lif !C)) \lif ((!A \lif !B) \lif (!A \lif !C))
\]
Kita bisa nggayuh kondisional tersarang kang pungkasan kanthi
ngetrapake \MP{} kaping pindho. Iki maneh ateges kita kepengin conto
saka \olref[prp]{ax:lif2} kang $!A \lif !C$-ne yaiku $!D \lif !D$,
!!{formula} kang dadi ancas kita. Mula, $!A$ lan $!C$ mesthi padha
karo~$!D$. Kepiye carane milih~$!B$ supaya $!A \lif (!B \lif !C)$
lan $!A \lif !B$, yaiku ing kasus kita $!D \lif (!B \lif !D)$ lan
$!D \lif !B$, uga bisa !!{derivable}? Formula kapisan wis dadi conto
saka \olref[prp]{ax:lif1}, apa wae kang dipilih dadi $!B$. Lan
$!D \lif !B$ bakal dadi conto liyane saka \olref[prp]{ax:lif1} yen
$!B$ yaiku $(!D \lif !D)$. Dadi, !!{derivation} kita yaiku:
\begin{derivation}
  1. & $!D \lif ((!D \lif !D) \lif !D)$ & \olref[prp]{ax:lif1}\\
  2. & $(!D \lif ((!D \lif !D) \lif !D)) \lif {}$\\
  & \qquad $((!D \lif (!D \lif !D)) \lif (!D \lif !D))$ & \olref[prp]{ax:lif2}\\
  3. & $(!D \lif (!D \lif !D)) \lif (!D \lif !D)$ & 1, 2, \MP\\
  4. & $!D \lif (!D \lif !D)$ & \olref[prp]{ax:lif1}\\
  5. & $!D \lif !D$ & 3, 4, \MP
\end{derivation}
\end{ex}

\begin{ex}\ollabel{ex:chain}
Kadhangkala kita arep nuduhake yen ana !!a{derivation} saka sawijining
!!{formula} saka !!{formula}s liyane~$\Gamma$. Upamane, ayo nuduhake
yen kita bisa !!{derive} $!A \lif !C$ saka $\Gamma = \{!A \lif !B,
!B \lif !C\}$.
\begin{derivation}
  1. & $!A \lif !B$ & \Hyp\\
  2. & $!B \lif !C$ & \Hyp\\
  3. & $(!B \lif !C) \lif (!A \lif (!B \lif !C))$ & \olref[prp]{ax:lif1} \\
  4. & $!A \lif (!B \lif !C)$ & 2, 3, \MP\\
  5. & $(!A \lif (!B \lif !C)) \lif {}$\\
  & \qquad $((!A \lif !B) \lif (!A \lif !C))$ & \olref[prp]{ax:lif2}\\
  6. & $((!A \lif !B) \lif (!A \lif !C))$ & 4, 5, \MP\\
  7. & $!A \lif !C$ & 1, 6, \MP
\end{derivation}
Larik-larik kanthi label ``\Hyp'' (kanggo ``hipotesis'') nuduhake yen
!!{formula} ing larik kasebut iku !!a{element} saka~$\Gamma$.
\end{ex}

\begin{prop}
  \ollabel{prop:chain} Yen $\Gamma \Proves !A \lif !B$ lan $\Gamma
  \Proves !B \lif !C$, mula $\Gamma \Proves !A \lif !C$
\end{prop}

\begin{proof}
  Upamana $\Gamma \Proves !A \lif !B$ lan $\Gamma \Proves !B \lif
  !C$. Banjur ana !!a{derivation} saka $!A \lif !B$ saka~$\Gamma$,
  lan uga ana !!a{derivation} saka~$!B \lif !C$ saka~$\Gamma$.
  Gabungna dadi siji !!{derivation} kanthi nggandhengake loro rerangken
  kasebut. Saiki tambahana larik 3--7 saka !!{derivation} ing conto
  sadurunge. Iki !!a{derivation} saka $!A \lif !C$---yaiku larik
  pungkasan ing !!{derivation} anyar---saka~$\Gamma$. Gatekna yen
  kabeneran larik 4 lan~7 tetep lumaku yen rujukan marang nomer
  larik~1 diganti nganggo rujukan marang larik pungkasan ing
  !!{derivation} saka~$!A \lif !B$, lan rujukan marang nomer larik~2
  diganti nganggo rujukan marang larik pungkasan ing !!{derivation}
  saka~$!B \lif !C$.
\end{proof}

\begin{prob} 
Nuduhna yen pratelan-pratelan iki lumaku kanthi menehi !!{derivation}s
saka aksioma:
\begin{enumerate} 
\item $(!A \land !B) \lif (!B \land !A)$
\item $((!A \land !B) \lif !C) \lif (!A \lif (!B \lif !C))$
\item $\lnot(!A \lor !B) \lif \lnot !A$
\end{enumerate} 
\end{prob}



\end{document}
